% ===================== RESUMEN EJECUTIVO (ABSTRACT ESTRUCTURADO) =====================
% Estilo Elsevier: párrafos con rótulos en negrita, sin viñetas.
% Extensión sugerida: máximo 250 palabras en total.
\newpage
\section*{Resumen Ejecutivo}
\addcontentsline{toc}{section}{Resumen Ejecutivo}

\textbf{Objetivo.}
% Indicar en una o dos frases qué se diseña, evalúa u optimiza, y bajo qué
% condiciones de operación o requisito del cliente.
[Describir el objetivo principal del proyecto y la necesidad que lo origina.]

\textbf{Métodos.}
% Citar el método de cálculo, correlaciones y normas aplicadas con su edición
% (p. ej. ASME VIII Div.1 ed. 2023), y las herramientas utilizadas.
[Describir metodología, correlaciones empleadas, normas aplicables y herramientas de cálculo.]

\textbf{Resultados.}
% Reportar solo los valores clave con unidades SI y su margen de diseño;
% remitir a las tablas y figuras del cuerpo del informe.
[Presentar los resultados numéricos principales con unidades y su trazabilidad.]

\textbf{Conclusiones.}
% Cerrar con la viabilidad técnica, la recomendación principal y, si aplica,
% la limitación más relevante del estudio.
[Enunciar las conclusiones y la recomendación principal del estudio.]

\vspace{1em}
\noindent\textbf{Palabras clave:}
% De cinco a ocho términos, separados por punto y coma, en español técnico.
[palabra clave 1]; [palabra clave 2]; [palabra clave 3]; [palabra clave 4]; [palabra clave 5].

\newpage
